\section{Introduction}

\IEEEPARstart{M}{otion} estimates from cameras with an endoscopic field of
view support navigation, reconstruction, and spatial registration. Their
evaluation, however, presupposes a reference trajectory. Robotic datasets use
forward kinematics~\cite{hayoz2023pose}; other laboratory systems use optical
trackers~\cite{hartwig2022constrained}. Neither is normally
available for a compact handheld camera.

A camera-mounted MEMS IMU is attractive because it is small, inexpensive,
and operates at a higher rate than the camera. Direct integration is
ill-conditioned: gravity, bias, scale factor, axis misalignment, timing, and
initial velocity all enter the position estimate. Data-driven inertial
odometry~\cite{chen2018ionet,yan2020ronin,liu2020tlio,zeinali2024imunet} can
learn motion priors, while online calibration such as
DUET~\cite{liu2023duet} learns sensor corrections. Yet a learned output is a
measurement only if its reference, calibration chain, uncertainty,
out-of-distribution behavior, and distance from the physical information
limit are reported.

Four coupled gaps motivate this work. First, a planar board may reproject
accurately while its tilt and translation remain strongly coupled. Second,
using the same IMU to ``improve'' its training labels creates information
leakage. Third, learned inertial odometry is normally reported as an accuracy
ranking between architectures, which hides how much of the residual is
physically unobservable at the chosen window and therefore common to every
architecture; without that decomposition it is impossible to tell whether a
better network or a better acquisition is required. Fourth, a simulated IMU
can cover motion that is unsafe or tedious to capture, but a digital twin is
useful only when its parameters have measured provenance and its domain gap is
quantified. These issues lead to the following research questions:
\begin{itemize}
  \item[RQ1)] How accurate and repeatable is a planar camera-pose reference
  under the spatial and temporal conditions of the instrument?
  \item[RQ2)] What protocol-dependent observability floor arises in windowed inertial
  displacement measurement for this instrument, and how much of the residual
  error is physically unobservable rather than unlearned?
  \item[RQ3)] Which acquisition property changes that floor, and by how much?
  \item[RQ4)] Which input/output representation approaches the observed floor, at what
  data and complexity cost, and does an audited digital twin reduce that
  cost?
\end{itemize}

\textbf{The contributions are:}
\begin{itemize}
  \item We establish a leakage-free camera--IMU measurement chain. We audit
  the planar pose reference, quantify its tilt--translation coupling, test
  camera--IMU rigidity on held-out increments, and show experimentally how
  IMU-corrected labels inflate the reported accuracy.
  \item We analyse the observability of windowed inertial displacement. The
  decomposition identifies an unobservable initial velocity, and a measured
  low-frequency velocity proxy explains the common empirical error floor of
  more than forty estimator variants, none of whose conventional architecture or augmentation changes improved consistently beyond the seed variation. A boundary experiment then quantifies
  the interval over which a confirmed pause improves inertial integration.
  \item We propose PhysNet, which combines analytic anchor-frame
  pre-integration with a densely supervised velocity-to-displacement network,
  a stillness-conditioned velocity gate, a learned lever arm and calibrated
  per-prediction uncertainty. It approaches the measured floor with one seventh
  of the parameters of adapted IMUNet.
  \item We build a parameter-provenance digital twin and a protocol-matched
  evaluation spanning synthetic test and OOD, direct sim-to-real transfer,
  session-wise cross-validation and a held-out real test. The study turns the
  pause rate and attitude quality into actionable acquisition parameters
  instead of attributing an empirical error floor to network architecture.
\end{itemize}

Fig.~\ref{fig:framework} summarizes the information flow and the separation
between reference, calibration, learning and held-out evaluation. The analysis
first determines which part of the measurand is observable from the sensor and
which part requires an acquisition constraint; the network comparison is then
interpreted relative to that limit.
